% ===== §18 =====
\section{The Remainder that Sustains the Bond}
\label{p27:sec:remainder}

\noindent
The paper has argued, across its length, for understanding, for the entry into the other's world, the crossing that generates a shared reality, the capacities by which the crossing is carried out. A reader who has followed this far might draw from it a conclusion the paper must now refuse, that the aim of relational understanding is a full understanding, a complete crossing in which the other's world is at last wholly entered and the two renderings brought into a total correspondence. This section argues, dialectically, that such a completion is neither possible nor to be wished, that the persistence of an intimate bond depends upon a remainder in the other that understanding does not close, and it draws from this a corrected conception of what relational understanding, rightly aimed, seeks.

The thesis is the paper's own argument to this point. Understanding is good, its work in the bond is real, and the capacities that carry it are worth cultivating. Left there, the argument inclines toward a horizon at which understanding would be complete, the other fully known, the two worlds mapped onto one another without remainder. This horizon is the same fusion the paper set out, in its first pages, to refuse, returning now in a subtler dress. It no longer promises the transparent transport of a message, and it promises instead a complete mutual observation, a bond in which each has entered the other so fully that nothing of the other remains unobserved. The subtler form is the more seductive, since it wears the appearance of the very understanding the paper commends, and it must be met on its own ground.

The antithesis meets it there. A full crossing of the other's world, were it reached, would end the bond it was meant to perfect, and the reason lies in the core introduced in the background. At the center of the other is an unsymbolizable remainder around which, on the psychoanalytic account, the other's desire is organized. Desire moves because it circles a core that no object and no observation fills, and the movement is the life of the desiring subject. Were the other's core at last reached, observed, brought into a complete correspondence with one's own rendering, the circling would have nothing left to circle. The remainder that drew desire would be closed, and with it the movement that the remainder sustained. A complete understanding of the other would not crown the bond with a final intimacy. It would still the desire that moved between the two, for it would fill the lack around which that desire turned, and a desire whose lack is filled is a desire extinguished. The completion the thesis reaches toward is, at its horizon, the death of the very movement that makes the bond alive.

This is not a defect that a better completion would avoid. It is the structure of desire, and it holds against any full crossing whatever. The point can be put in the paper's own terms, without the psychoanalytic vocabulary, so that it does not rest on a doctrine the reader might decline. The bond has been shown to live by generation, by the continued production, in the crossing of two worlds, of observations and values that neither world held before. Generation requires that there be, at each moment, something not yet observed from which the next observation is generated. A remainder in the other, a region of the other's world not yet crossed, is the reservoir from which the bond's further generation is drawn. Were the other fully crossed, this reservoir would be exhausted, and the generation that the bond lives by would have nothing further to draw upon. The bond would pass from a spiral that returns to the other and finds more, to a circle that returns to the other and finds the same, and the circle, for all its completeness, would be the end of the bond's growth. The remainder is not an obstacle to the bond's life. It is the condition of it.

There is a further point, of a normative kind, that joins the argument from desire and the argument from generation. To hold that one has fully understood the other, fully crossed their world, fully observed their core, is to treat the other as an object that a sufficient observation could exhaust. The other's unsymbolizable remainder is the mark of the other's standing as a subject and not an object, as a source of renderings that is never used up in one's observation of it. To reach, or to believe one has reached, a complete understanding of the other is thus to have annulled, in one's own regard, the other's subjectivity, to have converted a co-generating subject into a mapped and finished thing. The aspiration to full understanding, however tender its motive, meets at its limit the same wrong the normative case named, the reduction of the other to what one can possess of them. The remainder that understanding does not close is what keeps the other, in one's understanding of them, a subject who exceeds it.

\begin{claim}[The necessary remainder]
The aim of relational understanding is not a full understanding, and a complete crossing of the other's world would end the bond it was meant to perfect. The other's unsymbolizable remainder is the core around which desire moves, the reservoir from which the bond's further generation is drawn, and the mark of the other's standing as a subject who exceeds any observation of them. Were it closed, desire would be stilled, generation would be exhausted, and the other would be reduced to a finished object. The remainder that understanding does not reach is therefore the condition of the bond's persistence, and understanding rightly aimed does not seek to close it.
\end{claim}

The synthesis is not a midpoint between understanding and its refusal, and it does not counsel a strategic withholding, a deliberate leaving of the other unknown for the sake of preserving mystery, which would be one more possession, a management of the other's opacity for an effect one wants. It is a corrected conception of what understanding seeks. Relational understanding, rightly aimed, is not an asymptotic approach to a full crossing that it never quite completes. It is a sustained crossing that generates, at each meeting, a shared observation the two did not hold before, and that in the same motion honors the remainder from which the next crossing will be generated. The two moments are one. To understand the other well is to observe with the other what their crossing generates and to leave intact the core from which further generation comes, to enter without annexing, to observe without exhausting. This is the conception the normative case named as non-possession, now shown to have a further ground in the structure of desire and the condition of generation, and it is the conception the account has meant, throughout, by a crossing that observes with the other and does not possess them.

This synthesis does not resolve the tension it names, and it should not, for the tension is the living form of the bond. There is a real pull toward understanding, and it is right, and there is a real remainder that understanding must not close, and it is also right, and the bond lives in the holding of both, in a crossing that seeks the other and honors what in the other it will not reach. To resolve the tension in favor of full understanding would still the bond in a completed knowledge; to resolve it in favor of the remainder would abandon the crossing altogether and leave the two worlds unmet. The bond lives in neither resolution but in the sustained movement between them, the seeking that generates and the honoring that leaves room to generate again. The account holds the tension open, in keeping with its refusal of a final synthesis, since a bond, like the understanding that crosses it, is a movement to be sustained and not a matter to be finished.
